% defer/refcnt.tex
% mainfile: ../perfbook.tex
% SPDX-License-Identifier: CC-BY-SA-3.0

\section{引用计数}
\label{sec:defer:Reference Counting}
%
\epigraph{我永远不会放你走！}{佚名}

\IXalt{Reference counting}{reference count}
通过跟踪对给定对象的引用数量来防止该对象被过早释放。
因此，它有着悠久而光荣的使用历史，至少可以追溯到
1960 年代初 Weizenbaum 的一篇论文~\cite{Weizenbaum:1963:SLP:367593.367617}。
Weizenbaum 在讨论 reference counting 时把它当作已广为人知的概念，
因此它可能可以追溯到 1950 年代甚至 1940 年代。
而且可能更早，因为修理大型危险机器的工人长期以来一直使用
通过挂锁实现的机械式引用计数技术。
在进入机器之前，每个工人将一把挂锁锁在机器的开关上，
从而防止机器在该工人在内部时被通电。
因此，reference counting 是 Pre-BSD 路由并发实现的优秀且
历经考验的候选方案。

\begin{listing}
\input{CodeSamples/defer/route_refcnt@lookup.fcv}
\caption{Reference-Counted Pre-BSD 路由表查找（有缺陷！！！）}
\label{lst:defer:Reference-Counted Pre-BSD Routing Table Lookup}
\end{listing}

\begin{listing}
\input{CodeSamples/defer/route_refcnt@add_del.fcv}
\caption{Reference-Counted Pre-BSD 路由表添加\slash 删除（有缺陷！！！）}
\label{lst:defer:Reference-Counted Pre-BSD Routing Table Add/Delete}
\end{listing}

为此，
\cref{lst:defer:Reference-Counted Pre-BSD Routing Table Lookup}
展示了数据结构和 \co{route_lookup()} 函数，
\cref{lst:defer:Reference-Counted Pre-BSD Routing Table Add/Delete}
展示了 \co{route_add()} 和 \co{route_del()} 函数
（均在 \path{route_refcnt.c} 中）。
由于这些算法与\cref{lst:defer:Sequential Pre-BSD Routing Table}
中展示的顺序算法非常相似，因此只讨论差异之处。

\begin{fcvref}[ln:defer:route_refcnt:lookup:entry]
从\cref{lst:defer:Reference-Counted Pre-BSD Routing Table Lookup}开始，
\clnref{refcnt}添加了实际的引用计数器，
\clnref{freed}添加了 \co{->re_freed} 释放后使用检查字段，
\clnref{routelock}添加了用于同步并发更新的 \co{routelock}，
\end{fcvref}
\begin{fcvref}[ln:defer:route_refcnt:lookup:re_free]
\clnrefrange{b}{e}添加了 \co{re_free()}，它设置
\co{->re_freed}，使 \co{route_lookup()} 能够检查
释放后使用的错误。
\end{fcvref}
\begin{fcvref}[ln:defer:route_refcnt:lookup:lookup]
在 \co{route_lookup()} 本身中，
\clnrefrange{relprev:b}{relprev:e}释放前一个元素的引用计数，
如果计数变为零则释放它，
\clnrefrange{acq:b}{acq:e}获取新元素的引用，
\clnref{check_uaf,abort}执行释放后使用检查。
\end{fcvref}

\QuickQuiz{
	为什么要费心进行释放后使用检查？
}\QuickQuizAnswer{
	为了大大增加发现错误的概率。
	一个只分配和释放一种类型结构的小型压力测试程序
	（\path{routetorture.h}）可以容忍数量惊人的
	释放后使用行为。
	参见\cref{fig:debugging:Number of Tests Required for 99 Percent Confidence Given Failure Rate}
	（\cpageref{fig:debugging:Number of Tests Required for 99 Percent Confidence Given Failure Rate}）
	以及从\cpageref{sec:debugging:Hunting Heisenbugs}开始的
	\cref{sec:debugging:Hunting Heisenbugs}中的相关讨论，
	了解更多关于增加错误发现概率的重要性。
}\QuickQuizEnd

\begin{fcvref}[ln:defer:route_refcnt:add_del]
在\cref{lst:defer:Reference-Counted Pre-BSD Routing Table Add/Delete}中，
\clnref{acq1,rel1,acq2,rel2,rel3}引入了锁来同步并发更新。
\Clnref{init:freed}初始化 \co{->re_freed} 释放后使用检查字段，
最后 \clnrefrange{re_free:b}{re_free:e}在引用计数的新值为零时
调用 \co{re_free()}。
\end{fcvref}

\QuickQuiz{
	为什么\cref{lst:defer:Reference-Counted Pre-BSD Routing Table Add/Delete}中的
	\co{route_del()} 不使用引用计数来
	保护遍历到要释放的元素？
}\QuickQuizAnswer{
	因为遍历已经受到锁的保护，
	所以不需要额外的保护。
}\QuickQuizEnd

\begin{figure}
\centering
\resizebox{2.5in}{!}{\includegraphics{CodeSamples/defer/data/hps.2019.12.17a/perf-refcnt}}
\caption{受 Reference Counting 保护的 Pre-BSD 路由表}
\label{fig:defer:Pre-BSD Routing Table Protected by Reference Counting}
\end{figure}

\Cref{fig:defer:Pre-BSD Routing Table Protected by Reference Counting}
显示了 reference counting 在只读工作负载下的性能和可扩展性，
该工作负载使用十元素列表，运行在一个八路28核超线程 2.1\,GHz
x86 系统上，共有 448 个硬件线程
（\path{hps.2019.12.02a/lscpu.hps}）。
``ideal'' 曲线是通过运行\cref{lst:defer:Sequential Pre-BSD Routing Table}
中展示的顺序代码生成的，这之所以可行是因为这是一个只读工作负载。
reference counting 的性能极差，其可扩展性更是糟糕，
``refcnt'' 曲线与 x 轴无法区分。
鉴于\cref{chp:Hardware and its Habits}的内容，这不足为奇：
引用计数的获取和释放在原本只读的工作负载中增加了频繁的
共享内存写操作，从而招致了物理定律的严厉惩罚。
这完全是理所当然的，因为世界上所有的一厢情愿都不会
增加光速或减小现代数字电子产品中所用原子的大小。

\QuickQuizSeries{%
\QuickQuizB{
	为什么\cref{fig:defer:Pre-BSD Routing Table Protected by Reference Counting}
	中的``ideal''线在 224 个 CPU 处有一个断裂？
	它不应该是一条直线吗？
}\QuickQuizAnswerB{
	这个断裂是由超线程引起的。
	在这个特定系统中，每个插座中每个核心的第一个硬件线程
	有连续的 CPU 编号，
	接着是其他插座中每个核心的第一个硬件线程，
	最后是所有插座中每个核心的第二个硬件线程。
	在这个特定系统中，CPU 编号 0--27 是第一个插座的
	28 个核心中的第一个硬件线程，
	编号 28--55 是第二个插座的 28 个核心中的第一个硬件线程，
	依此类推，编号 196--223 是第八个插座的 28 个核心中的
	第一个硬件线程。
	然后 CPU 编号 224--251 是第一个插座的 28 个核心的
	第二个硬件线程，编号 252--279 是第二个插座的 28 个核心的
	第二个硬件线程，依此类推，直到编号 420--447 是第八个
	插座的 28 个核心的第二个硬件线程。

	这为什么重要？

	因为给定核心的两个硬件线程共享资源，
	而这个工作负载似乎允许单个硬件线程消耗其核心中
	超过一半的相关资源。
	因此，添加该核心的第二个硬件线程带来的增益比
	人们所期望的要少。
	其他工作负载可能从每个核心的第二个硬件线程中获得更大的收益，
	但这在很大程度上取决于硬件和工作负载的细节。
}\QuickQuizEndB
%
\QuickQuizE{
	\cref{fig:defer:Pre-BSD Routing Table Protected by Reference Counting}
	中的 refcnt 曲线难道不应该至少稍微偏离 x 轴吗？
}\QuickQuizAnswerE{
	请定义``稍微''。

\begin{figure}
\centering
\resizebox{2.5in}{!}{\includegraphics{CodeSamples/defer/data/hps.2019.12.17a/perf-refcnt-logscale}}
\caption{受 Reference Counting 保护的 Pre-BSD 路由表，对数刻度}
\label{fig:defer:Pre-BSD Routing Table Protected by Reference Counting; Log Scale}
\end{figure}

	\Cref{fig:defer:Pre-BSD Routing Table Protected by Reference Counting; Log Scale}
	显示了相同的数据，但使用对数-对数图。
	如你所见，refcnt 曲线在两个 CPU 时降到了 5,000 以下。
	这意味着两个 CPU 时的 refcnt 性能比
	\cref{fig:defer:Pre-BSD Routing Table Protected by Reference Counting}
	中第一个 y 轴刻度 $5 \times 10^6$ 小一千倍以上。
	因此，\cref{fig:defer:Pre-BSD Routing Table Protected by Reference Counting}
	中展示的 reference counting 性能描述过于准确了。
}\QuickQuizEndE
}

但情况还会更糟。

多个更新线程反复调用 \co{route_add()} 和 \co{route_del()} 将很快
遇到\cref{lst:defer:Reference-Counted Pre-BSD Routing Table Lookup}
\clnrefr{ln:defer:route_refcnt:lookup:lookup:abort}上的 \co{abort()} 语句，
这表明存在释放后使用错误。
这反过来意味着引用计数不仅严重降低了可扩展性和性能，
而且未能提供所需的保护。

导致释放后使用错误的一个事件序列如下，
给定\cref{fig:defer:Pre-BSD Packet Routing List}中所示的列表：

\begin{fcvref}[ln:defer:route_refcnt:lookup]
\begin{enumerate}
\item	线程~A 查找地址~42，到达
	\cref{lst:defer:Reference-Counted Pre-BSD Routing Table Lookup}中
	\co{route_lookup()} 的\clnref{lookup:check_NULL}。
	换句话说，线程~A 有一个指向第一个元素的指针，
	但尚未获取对它的引用。
\item	线程~B 调用\cref{lst:defer:Reference-Counted Pre-BSD Routing Table Add/Delete}
	中的 \co{route_del()} 来删除地址~42 的路由条目。
	它成功完成，由于该条目的 \co{->re_refcnt} 字段等于值 1，
	它调用 \co{re_free()} 来设置 \co{->re_freed} 字段并释放该条目。
\item	线程~A 继续执行 \co{route_lookup()}。
	它的 \co{rep} 指针非 \co{NULL}，但
	\clnref{lookup:check_uaf}发现其 \co{->re_freed} 字段非零，
	因此\clnref{lookup:abort}调用 \co{abort()}。
\end{enumerate}
\end{fcvref}

问题在于引用计数位于要保护的对象中，但这意味着在获取引用计数本身的
那一瞬间没有保护！
这是 Gamsa 等人~\cite{Gamsa99}注意到的锁问题的 reference counting 对应物。
可以想象使用全局锁或引用计数来保护每个路由条目的引用计数获取，
但这会导致严重的竞争问题。
虽然存在允许在并发环境中安全获取引用计数的
算法~\cite{Valois95a}，但它们不仅极其复杂且容易出错~\cite{MagedMichael95a}，
而且提供的性能和可扩展性也很差~\cite{ThomasEHart2007a}。

简而言之，并发确实大大降低了 reference counting 的实用性！
当然，与其他同步原语一样，引用计数也有众所周知的易用性缺陷。
这些缺陷一方面可能导致内存泄漏，另一方面可能导致过早释放。

这就是在\cpageref{sec:defer:Mysteries reference-counting trap}中
提到的等待并发代码粗心开发者的 reference counting 陷阱。

\QuickQuiz{
	如果并发``确实大大降低了 reference counting 的实用性''，
	为什么 Linux 内核中有那么多引用计数器？
}\QuickQuizAnswer{
	那句话确实说的是``降低了实用性''，
	而不是``消除了实用性''，不是吗？

	请参见\cref{sec:together:Refurbish Reference Counting}，
	其中讨论了 Linux 内核在高度并发环境中
	利用 reference counting 的一些技术。
}\QuickQuizEnd

有时候以完全不同的方式看待问题有助于成功解决它。
为此，下一节描述了一种可以被认为是内翻的引用计数，
它提供了不错的性能和可扩展性。
